% =======================================================
% 1. INTRODUZIONE
% =======================================================
\section{Introduzione}

\subsection{Scopo del Documento}
Il presente documento ha lo scopo di raccogliere e definire in modo univoco i termini
specialistici utilizzati all'interno della documentazione del gruppo \textit{Synergo Unit}.
La presenza di un glossario strutturato permette di evitare ambiguità terminologiche,
favorire una comunicazione chiara e garantire coerenza tra i diversi artefatti prodotti
nel corso del progetto di Ingegneria del Software A.A.~2025/2026.

Ogni termine incluso è stato selezionato in quanto rilevante per il dominio del progetto,
per il processo di sviluppo o per le convenzioni interne adottate dal gruppo.  
Il glossario costituisce pertanto un riferimento condiviso e aggiornato per tutti i membri
del team e per gli stakeholder esterni.

\subsection{Regole di Tipografia}
Per facilitare la lettura e l’individuazione dei termini definiti nel presente documento,
il gruppo adotta la seguente convenzione tipografica:

\begin{itemize}
    \item la \textbf{prima occorrenza} di un termine presente nel glossario è contrassegnata
    dal \textbf{pedice G}, ad esempio: \textit{termine\textsubscript{G}};
    \item le occorrenze successive dello stesso termine non riportano il pedice, al fine di
    evitare ridondanze e mantenere il testo più scorrevole;
    \item i termini sono organizzati in ordine alfabetico e presentati con una definizione
    chiara, concisa e non ambigua.
\end{itemize}

\subsection{Riferimenti}

\subsubsection{Normativi}

\begin{itemize}
    \item \link{https://www.math.unipd.it/~tullio/IS-1/2025/Progetto/C6p.pdf}{Capitolato C6 -- Second Brain (Zucchetti S.p.A.)};

    \item \link{https://www.math.unipd.it/~tullio/IS-1/2025/Dispense/PD1.pdf}{Regolamento del Progetto Didattico -- Corso di Ingegneria del Software, A.A. 2025/2026};

    \item \link{https://synergo-unit-unipd.github.io/Second-Brain/docs/RTB/doc_gestionali/doc_interni/norme_di_progetto/norme_di_progetto.pdf}{Norme di Progetto\textsubscript{G}, v2.12.0}.
\end{itemize}

\subsubsection{Informativi}

\begin{itemize}
    \item \link{https://synergo-unit-unipd.github.io/Second-Brain/docs/RTB/doc_tecnici/doc_esterni/analisi_dei_requisiti/analisi_dei_requisiti.pdf}{Analisi dei Requisiti\textsubscript{G}, v1.0.0};

    \item \link{https://synergo-unit-unipd.github.io/Second-Brain/docs/RTB/doc_gestionali/doc_esterni/piano_di_progetto/piano_di_progetto.pdf}{Piano di Progetto\textsubscript{G}, v1.18.1};

    \item \link{https://synergo-unit-unipd.github.io/Second-Brain/docs/RTB/doc_gestionali/doc_esterni/piano_di_qualifica/piano_di_qualifica.pdf}{Piano di Qualifica\textsubscript{G}, v1.3.0};

    \item \link{https://synergo-unit-unipd.github.io/Second-Brain/docs/RTB/doc_tecnici/doc_esterni/stato_dell_arte/stato_dell_arte.pdf}{Analisi dello Stato dell'Arte\textsubscript{G}, v1.0.0};

    \item Slide del corso di Ingegneria del Software, A.A. 2025/2026.
\end{itemize}
